\section{Related Work}
\label{sec:related}

Several smartphone applications have been developed to attempt to detect rogue
base stations. Park~\etal evaluated five popular IMSI catcher
applications---\emph{SnoopSnitch}~\cite{snoopsnitch}, \emph{Cell Spy Catcher},
\emph{GSM Spy Finder}, \emph{Darshak}~\cite{darshak}, and the \emph{Android IMSI
Catcher Detector (AIMSICD)}~\cite{AIMSICD},---for Android in order to evaluate
their efficacy against rogue base stations using a rogue base station framework
they developed called \emph{White-Stingray}~\cite{whitestingray}. Since this
work was published in 2017, several of these applications are no longer
available in the Google Play Store. 

More recently, Arnold~\etal, developed \emph{CellGuard}, an iOS application for
analyzing nearby cellular networks to detect possible rogue base
stations~\cite{cellguard}. \emph{CellGuard} uses several factors, including
whether or not a base station is present in Apple's \cps, to make this
determination. Unlike their work, which focuses on detecting a local adversary
in real time, this work seeks to map the world's cellular infrastructure and
detect malicious and suspicious rogue base stations longitudinally across time
and space.  

Most closely related to this work is Rye and Levin's investigation of the
security and privacy vulnerabilities associated with Apple's
\ac{WPS}~\cite{rye2024surveilling}. The authors demonstrated that an
unprivileged attacker can amass hundreds of millions of \wifi \ac{BSSID}
geolocations by querying Apple's \ac{WPS} with little \emph{a priori} knowledge.
In this work, we 
